\documentclass[12pt,a4paper]{report}
\usepackage[utf8]{inputenc}
\usepackage{geometry}
\geometry{margin=1in}
\usepackage{listings}
\usepackage{xcolor}
\usepackage{hyperref}
\usepackage{titlesec}
\usepackage{graphicx}
\usepackage{array}

% Colors for code highlighting
\definecolor{codegreen}{rgb}{0,0.6,0}
\definecolor{codegray}{rgb}{0.5,0.5,0.5}
\definecolor{codepurple}{rgb}{0.58,0,0.82}
\definecolor{backcolour}{rgb}{0.95,0.95,0.92}

\lstdefinestyle{mystyle}{
    backgroundcolor=\color{backcolour},   
    commentstyle=\color{codegreen},
    keywordstyle=\color{magenta},
    numberstyle=\tiny\color{codegray},
    stringstyle=\color{codepurple},
    basicstyle=\ttfamily\footnotesize,
    breakatwhitespace=false,         
    breaklines=true,                 
    captionpos=b,                    
    keepspaces=true,                 
    numbers=left,                    
    numbersep=5pt,                  
    showspaces=false,                
    showstringspaces=false,
    showtabs=false,                  
    tabsize=2
}

\lstset{style=mystyle}

\title{
    \Huge \textbf{Zoho OS} \\
    \large \textit{Building a 64-bit Operating System from Scratch} \\
    \vspace{1cm}
    \normalsize Technical Documentation \& Architectural Blueprint
}
\author{
    \textbf{Atharva Bodade} \\
    \small Shri Sant Gajanan Maharaj College of Engineering, Shegaon \\
    \vspace{0.5cm}
    \small \url{https://github.com/UnsettledAverage73/zoho-os}
}
\date{\today}

\begin{document}

\maketitle

\tableofcontents
\newpage

\chapter{Project Introduction: The Core Philosophy}
\section{The Vision: From Silicon to Software}
When you turn on your computer, a miracle happens. In less than 10 seconds, a piece of silicon and metal turns into a workspace. This project, \textbf{Zoho OS}, is the journey of building that miracle from the ground up. We are not using pre-made kernels, third-party libraries, or existing operating systems as a base; we are communicating directly with the hardware at the most fundamental level.

The source code for this project is available at: \\
\url{https://github.com/UnsettledAverage73/zoho-os}

The philosophy behind Zoho OS is \textit{"Total Understanding."} By writing every line of code—from the bootloader handshake to the graphical user interface—we eliminate the "black boxes" of modern computing and reclaim control over the machine.

\section{The Core Objectives}
The objective is to create a functional, high-performance, 64-bit monolithic operating system. The project focuses on several key technical pillars:
\begin{itemize}
    \item \textbf{Bare-Metal Execution:} Running directly on the hardware without any underlying abstraction layer.
    \item \textbf{Modern Graphics Architecture:} Implementing a high-resolution linear framebuffer to support millions of colors and smooth animations.
    \item \textbf{Symmetric Multi-Processing (SMP):} Utilizing all available CPU cores to enable true parallel execution.
    \item \textbf{Advanced Memory Management:} Implementing 4-level paging to provide every process with a secure, 64-bit virtual address space.
    \item \textbf{Custom System Logic:} Building our own interrupt handlers, task schedulers, and I/O drivers.
\end{itemize}

\chapter{Architecture \& Design}
\section{The Kitchen Analogy}
To understand the kernel, we use the analogy of a restaurant:
\begin{itemize}
    \item \textbf{The Hardware (Kitchen):} The stove and ingredients. It has power but no direction.
    \item \textbf{The Apps (Customers):} They want results but don't know how to use the hardware.
    \item \textbf{The Kernel (The Chef):} This is Zoho OS. It manages the hardware resources to serve the applications safely and efficiently.
\end{itemize}

\section{The Folder Structure}
The project is organized to separate the low-level "starting" code from the high-level "operating" code.
\begin{itemize}
    \item \textbf{src/boot/:} The "First Responders." Contains .asm files that wake up the CPU and transition to 64-bit mode.
    \item \textbf{src/kernel/:} The "Engine Room." Contains .c and .asm files for core logic: memory management, task switching, and hardware drivers.
    \item \textbf{include/kernel/:} The "Blueprint." Contains .h files defining structures and function signatures.
    \item \textbf{isodir/:} The "Packaging." Temporary folder used to create the final .iso file.
    \item \textbf{Makefile:} The "Recipe" that tells the computer how to compile everything.
    \item \textbf{linker.ld:} The "Map" that tells the compiler where each piece of code lives in RAM.
\end{itemize}

\section{The Tech Stack}
\begin{itemize}
    \item \textbf{x86\_64 Assembly (NASM):} For code C cannot do (like talking directly to CPU registers).
    \item \textbf{C (GCC):} For 99\% of kernel logic. Used in "Freestanding" mode (no standard library).
    \item \textbf{Build System:} GNU Make. Automates the compilation process.
    \item \textbf{Boot Protocol:} Multiboot2. Standardized way for GRUB to hand over control.
    \item \textbf{Emulation:} QEMU. A virtual computer used to test the OS.
\end{itemize}

\chapter{Phase 1: Bootstrapping \& 64-bit Transition}
\section{The Multiboot2 Header (src/boot/header.asm)}
When a computer turns on, it runs firmware (BIOS/UEFI) which looks for a Bootloader (like GRUB). GRUB looks for a specific signature called a Multiboot2 Header.

\begin{lstlisting}[language={[x86masm]Assembler}, caption=src/boot/header.asm]
section .multiboot_header
header_start:
    ; Magic number (multiboot 2)
    dd 0xe85250d6
    ; Architecture 0 (protected mode i386)
    dd 0
    ; Header length
    dd header_end - header_start
    ; Checksum
    dd 0x100000000 - (0xe85250d6 + 0 + (header_end - header_start))

    ; Framebuffer tag request
    align 8
    dw 5    ; type
    dw 0    ; flags
    dd 20   ; size
    dd 1024 ; width
    dd 768  ; height
    dd 32   ; depth

    ; End tag
    align 8
    dw 0
    dw 0
    dd 8
header_end:
\end{lstlisting}

\subsection{Line-by-Line Breakdown}
\begin{itemize}
    \item \textbf{section .multiboot\_header:} Tells the linker to put this at the very beginning of the binary.
    \item \textbf{dd 0xe85250d6:} The Magic Number for Multiboot2.
    \item \textbf{dd 0:} Architecture flag (32-bit protected mode).
    \item \textbf{dd 0x100000000 - (...):} Checksum to verify the header is not corrupted.
    \item \textbf{Framebuffer Tag:} Requests GRUB to set up a 1024x768 graphics mode.
\end{itemize}

\section{The 32-bit Entry \& Environment Checks (src/boot/main.asm)}
GRUB jumps to the label \textit{start}. We are now in 32-bit Protected Mode.

\begin{lstlisting}[language={[x86masm]Assembler}, caption=src/boot/main.asm]
global start
extern long_mode_start
section .text
bits 32
start:
    cli                 ; 1. Disable interrupts
    mov esp, stack_top  ; 2. Set up a stack
    
    ; Preserve multiboot info early!
    mov edi, eax        ; Magic number
    mov esi, ebx        ; Pointer to multiboot info structure

    call check_cpuid
    call check_long_mode
    call setup_page_tables
    call enable_paging
    lgdt [gdt64.pointer]
    jmp gdt64.code:long_mode_start
\end{lstlisting}

\subsection{Breakdown}
\begin{itemize}
    \item \textbf{cli:} Disables interrupts as we don't have an IDT yet.
    \item \textbf{mov esp, stack\_top:} Sets up a 16KB stack for C functions.
    \item \textbf{Preserving eax and ebx:} Saves Multiboot info passed by GRUB into registers for later use in C.
    \item \textbf{check\_cpuid / check\_long\_mode:} Verifies the CPU supports 64-bit mode.
    \item \textbf{setup\_page\_tables:} Creates identity-mapped page tables.
    \item \textbf{enable\_paging:} Flips bits in CR0, CR4, and EFER to enable Long Mode.
    \item \textbf{jmp gdt64.code:long\_mode\_start:} A "Far Jump" to switch the CPU to 64-bit mode.
\end{itemize}

\section{Transition to 64-bit (src/boot/long\_mode\_init.asm)}
This is the first file where we use 64-bit instructions.

\begin{lstlisting}[language={[x86masm]Assembler}, caption=src/boot/long\_mode\_init.asm]
global long_mode_start
extern kmain
section .text
bits 64
long_mode_start:
    ; 1. Nullify Data Segment Registers
    mov ax, 0
    mov ss, ax
    mov ds, ax
    mov es, ax
    mov fs, ax
    mov gs, ax

    ; 2. Print 'OK' to VGA
    mov rax, 0x0f4b0f4f
    mov [0xb8000], rax

    ; 3. The Big Leap
    call kmain
    hlt
\end{lstlisting}

\chapter{Phase 2: Kernel Initialization}
Welcome to the brain of the operating system. We have successfully crossed from Assembly into the structured world of C.

\section{kmain() - The Orchestrator}
In \textit{src/kernel/main.c}, \textit{kmain()} wakes up every part of the computer in the correct order.

\begin{lstlisting}[language=C, caption=src/kernel/main.c]
void kmain(unsigned long magic, unsigned long addr) {
    struct multiboot_info* mb_info = (struct multiboot_info*)addr;
    
    cpu_early_init();
    serial_init();
    vga_clear();
    klog_set_level(LOG_DEBUG);
    klog(LOG_INFO, "KERNEL", "Zoho OS Booting...");

    pmm_init(mb_info);
    vmm_init();
    kmalloc_init();

    gdt_init();
    idt_init_global();
    idt_init_per_cpu();

    __asm__ volatile ("sti"); // Enable Interrupts
    
    while(1) {
        serial_poll_input();
        e1000_poll();
        window_update();
        window_draw_all();
    }
}
\end{lstlisting}

\chapter{Phase 3: Hardware Interrupts \& IDT}
An Interrupt is like a phone call. It can be a hardware event (key press) or a software emergency (division by zero).

\section{The Interrupt Descriptor Table (IDT)}
The IDT acts as a "Phone Directory" with 256 entries. Each entry tells the CPU where to jump when an interrupt happens.

\begin{lstlisting}[language=C, caption=include/kernel/idt.h]
struct idt_entry {
    uint16_t offset_low;    // Lower 16 bits of handler address
    uint16_t selector;      // Code segment selector (GDT)
    uint8_t ist;            // Interrupt Stack Table
    uint8_t type_attr;      // Flags (present, ring level)
    uint16_t offset_middle; // Middle 16 bits
    uint32_t offset_high;   // Upper 32 bits
    uint32_t reserved;
} __attribute__((packed));
\end{lstlisting}

\chapter{Phase 4: Memory Management}
Memory management is divided into three layers:
\begin{enumerate}
    \item \textbf{Physical Memory Management (PMM):} Managing physical RAM frames (4KB chunks).
    \item \textbf{Virtual Memory Management (VMM):} Creating virtual addresses and 4-level page tables (PML4, PDPT, PD, PT).
    \item \textbf{Kernel Heap (kmalloc):} Providing dynamic allocation for small chunks of memory.
\end{enumerate}

\chapter{Phase 5: Multitasking \& Multiprocessing}
\section{Symmetric Multi-Processing (SMP)}
We use ACPI MADT tables to find all CPU cores and wake them up using a "Trampoline" code (started in 16-bit real mode) and IPI (Inter-Processor Interrupts).

\section{Task Management}
Every program is represented by a \textit{task\_t} struct (Process Control Block). Context switching involves saving the current CPU state and loading the state of the next task by swapping the \textit{RSP} (Stack Pointer) register.

\chapter{The Build System and Packaging}
The Makefile automates the entire process:
\begin{enumerate}
    \item \textbf{Compiling:} Converts .c and .asm files into .o object files.
    \item \textbf{Linking:} Glues object files into \textit{kernel.bin} using \textit{linker.ld}.
    \item \textbf{ISO Creation:} Packages the kernel and GRUB config into a bootable .iso using \textit{grub-mkrescue}.
\end{enumerate}

\chapter{Summary of the Total Picture}
\begin{center}
\begin{tabular}{|m{3cm}|m{2.5cm}|m{8cm}|}
\hline
\textbf{Concept} & \textbf{File} & \textbf{Purpose} \\
\hline
System Calls & syscall.c & Lets user apps talk to the kernel safely. \\
\hline
PCI Scanning & pci.c & Finds hardware (NIC, Sound, Disk). \\
\hline
Linking & linker.ld & Arranges code in RAM correctly. \\
\hline
Shell & shell.c & The command-line interface for the user. \\
\hline
String Lib & string.c & Manual versions of memcpy, strlen, etc. \\
\hline
\end{tabular}
\end{center}

\chapter{Conclusion}
From the first line of Assembly to a fully functional GUI, we have explored the entire lifecycle of an Operating System. Zoho OS demonstrates the fundamental "magic" that happens every time a computer is powered on.

\end{document}
